\documentclass[oneside]{amsart}
\usepackage[all]{xy}
\usepackage{tikz-cd}
\usepackage[T1]{fontenc}
\usepackage{xstring}
\usepackage{xparse}
\usepackage{xr-hyper}
\usepackage{xcolor}
\definecolor{brightmaroon}{rgb}{0.76, 0.13, 0.28}
\usepackage[linktocpage=true,colorlinks=true,hyperindex,citecolor=blue,linkcolor=brightmaroon]{hyperref}
\usepackage[nameinlink]{cleveref}
\usepackage[left=1.25in,right=1.25in,top=0.75in,bottom=0.75in]{geometry}
\setlength{\marginparwidth}{2cm}
%\usepackage[charter,greekfamily=didot]{mathdesign}
%\usepackage{Baskervaldx}
\usepackage{amssymb}
\usepackage{stmaryrd}
\usepackage{mathrsfs}
\usepackage{mathpazo}
\usepackage{eulervm}
\linespread{1.05}

\usepackage[nobottomtitles]{titlesec}
\usepackage{marginnote}
\usepackage{enumerate}
\usepackage{longtable}
\usepackage{aurical}
\usepackage{microtype}
\usepackage{mathtools}

\newtheoremstyle{ega-env-style}%
{}{}{\rmfamily}{}{\bfseries}{.}{ }{\thmnote{(#3)}}%

\newtheoremstyle{ega-thm-env-style}%
{}{}{\itshape}{}{\bfseries}{. --- }{ }{\thmname{#1}\thmnumber{ #2}\thmnote{ (#3)}}%

\newtheoremstyle{ega-defn-env-style}%
{}{}{\rmfamily}{}{\bfseries}{. --- }{ }{\thmname{#1}\thmnumber{ #2}\thmnote{ (#3)}}%

\theoremstyle{ega-env-style}
\newtheorem*{env}{---}

\theoremstyle{ega-thm-env-style}
\newtheorem{theorem}{Theorem}[section]
\newtheorem{proposition}[theorem]{Proposition}
\newtheorem{lemma}[theorem]{Lemma}
\newtheorem{corollary}[theorem]{Corollary}
\newtheorem{conjecture}[theorem]{Conjecture}

\theoremstyle{ega-defn-env-style}
\newtheorem{definition}[theorem]{Definition}
\newtheorem{example}[theorem]{Example}
\newtheorem*{examples}{Examples}
\newtheorem*{remark}{Remark}
\newtheorem*{remarks}{Remarks}
\newtheorem*{notation}{Notation}
\newtheorem*{exercise}{Exercise}
\newtheorem*{properties}{Properties}
\newtheorem*{consequences}{Consequences}
\newtheorem*{note}{Note}
\newtheorem*{fact}{Fact}
\newtheorem*{claim}{Claim}

\makeatletter
\def\l@subsection{\@tocline{2}{0pt}{2.5pc}{2.2pc}{}}
\def
\section{\@startsection{section}{1}%
	\z@{.7\linespacing\@plus\linespacing}{.5\linespacing}%
{\normalfont\bfseries\Large\scshape\centering}}
\renewcommand{\@seccntformat}[1]{%
	\ifnum\pdfstrcmp{#1}{section}=0\textsection\fi%
\csname the#1\endcsname.~}
\makeatother

\def\mathcal{\mathscr}
%% Fonts

\def\sh{\mathcal}                   % sheaf font
\def\bb{\mathbf}                    % bold font
\def\cat{\mathsf}                   % category font
\def\numfont{\mathbb}

%% Font Letters

\def\CA{\mathcal{A}}
\def\CE{\mathcal{E}}
\def\CF{\mathcal{F}}
\def\CG{\mathcal{G}}
\def\CL{\mathcal{L}}
\def\OO{\mathcal{O}}
\def\CX{\mathcal{X}}
\def\b1{\mathbb{1}}

%% Cohomology

\def\CH{\mathrm{H}}                 % cohomology H
\def\CHH{\check{\HH}}               % Čech cohomology H
\def\RD{\mathrm{R}}                 % right derived R
\def\LD{\mathrm{L}}                 % left derived L
\def\dual#1{{#1}^\vee}              % dual
\def\Tor{\operatorname{Tor}}        % Tor
\def\Ext{\operatorname{Ext}}        % Ext
\def\HdR{\mathrm{H}_{\mathrm{dR}}}  % de Rham cohomology
\def\Zc{\underline{\mathbb{Z}}}     % constant sheaf with integer coeffs

%% Categories

\def\A{\cat{A}}                     % category A, usually abelian
\def\C{\cat{C}}                     % category C
\def\op{^\cat{op}}                  % opposite category
\def\Set{\cat{Sets}}                % category of sets
\def\Grp{\cat{Gps}}                 % category of groups
\def\Alg{\cat{Alg}}                 % category of algebras
\def\QCoh{\cat{QCoh}}               % category of quasicoherent sheaves
\def\supertilde{{\,\widetilde{\,}\,}}   % use \supertilde instead of ^\sim
\def\Map{\operatorname{Map}}        % morphisms (any category)
\def\Hom{\operatorname{Hom}}        % morphisms (additive category)
\def\iHom{\underline{\Hom}}         % interal hom
\def\End{\operatorname{End}}        % endomorphisms
\def\Aut{\operatorname{Aut}}        % automorphisms
\def\ker{\operatorname{ker}}        % kernel
\def\img{\operatorname{im}}         % image
\def\coker{\operatorname{coker}}    % cokernel
\def\pr{\operatorname{pr}}          % projection
\def\EssIm{\operatorname{EssIm}}    % essential image
\DeclareMathOperator*{\colim}{colim}   % colimit

\def\DAet{\cat{DA}\et}
\def\SH{\cat{SH}}
\def\Sh{\cat{Sh}}                   % category of sheaves
\def\PSh{\cat{PSh}}                 % category of presheaves

%% Schemes

\def\Proj{\operatorname{Proj}}      % Proj
\def\Supp{\operatorname{Supp}}      % support
\def\Spec{\operatorname{Spec}}      % Spec
\def\Spf{\operatorname{Spf}}        % formal Spec
\def\Aff{\mathbb A}                 % affine space
\def\P{\mathbb{P}}                  % projective space
\def\Pic{\operatorname{Pic}}        % Picard group

%% Standard Operators

\def\codim{\operatorname{codim}}    % codimension
\def\id{\operatorname{id}}          % identity

%% Arrows
\renewcommand{\to}{\mathchoice{\longrightarrow}{\rightarrow}{\rightarrow}{\rightarrow}}
\newcommand{\from}{\mathchoice{\longleftarrow}{\leftarrow}{\leftarrow}{\leftarrow}}
\let\mapstoo\mapsto
\renewcommand{\mapsto}{\mathchoice{\longmapsto}{\mapstoo}{\mapstoo}{\mapstoo}}
\def\isoto{\simeq}                  % isomorphism
\def\simto{\xrightarrow{\sim}}      % isomorphism arrow
\def\surjto{\twoheadrightarrow}     % surjetion
\def\injto{\hookrightarrow}         % injection

%% Under/Over Accents

\newcommand{\wh}[1]{\widehat{#1}}   % hat
\newcommand{\wt}[1]{\widetilde{#1}}    % tilde
\def\ul{\underline}                % underline

%% Spaces

\def\Z{\numfont Z}                  % integers
\def\Q{\numfont Q}                  % rationals
\def\N{\numfont N}                  % naturals
\def\R{\numfont R}                  % reals
\def\H{\mathcal H}

%% Groups

\def\GL{\bb{GL}}                   % general linear group
\def\SL{\bb{SL}}                   % special linear group
\def\Sp{\bb{Sp}}
\def\PSp{\bb{PSp}}
\def\GSp{\bb{GSp}}
\def\det{\operatorname{det}}       % determinant

%% Sub/Superscripts

\def\et{^\text{\'et}}              % \'etale
\def\an{^\text{an}}                % analytic
\def\th{^\text{th}}                % th

%% Misc

\newcommand{\defn}[1]{\textbf{#1}}  % definition highlighting
\def\defeq{:=}                     % definition equation
\def\eps{\varepsilon}              % correct epsilon
\newcommand{\gen}[1]{\left\langle\!\left\langle #1 \right\rangle\!\right\rangle}


\def\shHom{\sh{H}\textup{\kern-2.2pt{\Fontauri\slshape om}}}   % sheaf Hom
\def\shProj{\sh{P}\textup{\kern-2.2pt{\Fontauri\slshape roj}}} % sheaf Proj
\def\shExt{\sh{E}\textup{\kern-2.2pt{\Fontauri\slshape xt}}}   % sheaf Ext
\def\shGr{\sh{G}\textup{\kern-2.2pt{\Fontauri\slshape r}}}     % sheaf Gr
\def\shDer{\sh{D}\,\textup{\kern-2.2pt{\Fontauri\slshape er}}} % sheaf Der
\def\shDiff{\sh{D}\,\textup{\kern-2.2pt{\Fontauri\slshape if{}f}}\,} % sheaf Diff
\def\shHomcont{\sh{H}\textup{\kern-2.2pt{\Fontauri\slshape om.\,cont}}}   % sheaf Hom.cont
\def\shAut{\sh{A}\textup{\kern-2.2pt{\Fontauri\slshape ut}}}   % sheaf Aut
\def\shSym{\sh{S}\textup{\kern-2.2pt{\Fontauri\slshape ym}}}   % sheaf Sym

\makeatletter
\newcommand{\cbigoplus}{\DOTSB\cbigoplus@\slimits@}
\newcommand{\cbigoplus@}{\mathop{\widehat{\bigoplus}}}
\makeatother

\usepackage[backend=biber,style=alphabetic]{biblatex}
\addbibresource{pwg.bib}

\renewcommand{\A}{\mathbb{A}}
\newcommand{\Bred}{\mathcal{B}_\text{red}}
\newcommand{\tf}{\mathfrak{t}}
\newcommand{\muT}{\mu_{\mathrm{Tam}}}
\newcommand{\F}{\mathbb{F}}
\DeclareMathOperator{\Bun}{Bun}
\DeclareMathOperator{\Tr}{Tr}
\DeclareMathOperator{\Frob}{Frob}
\DeclareMathOperator{\Ran}{Ran}
\newcommand{\red}{_{\mathrm{red}}}
\newcommand{\km}{\mathfrak{m}}
\DeclareMathOperator{\Sym}{Sym}
\DeclareMathOperator{\Shv}{Shv}
\DeclareMathOperator{\Gr}{Gr}
\def\D{\mathbb{D}}

\title{Introduction}
\begin{document}


\maketitle

\section{Motivation}

We can begin our story with a very classical problem from number theory: classifying integral quadratic forms (even lattices) up to isomorphism.

This is in itself an hopelessly hard problem, but we can break it up by grouping lattices by ``genus''. $\Lambda,\Lambda'$ have the same genus if they become isomorphic over $\R$ and over $\Z_p$ for all prime $p$. Note that, even though the Hasse-Minkowski theorem tells us that isomorphism classes of quadratic forms over $\Q$ are determined by their base change to $\R$ and $\Q_p$ for all $p$, the integral situation is quite different. Indeed, there can be many (though only finitely many) different lattices of the same genus. It is therefore a natural question to ask how to count the number of lattices in a given genus. The main result in this direction is the Siegel Mass formula:

\begin{theorem}[Siegel]
Let $\Lambda$ be a positive definite lattice of rank $n$. Then we have a product formula
\[
\left| \left\{ \text{lattices of genus }\Lambda \right\}  \right| = c_n \prod_p m_p(\Lambda),
\] 
for some numbers $m_p(\Lambda)$ which depends only on the isomorphism class of $\Lambda\otimes \Z_p$ and numbers $c_n$ which depend only on $n$ (equivalently, on the isomorphism class of $\Lambda\otimes \R$).
\end{theorem}
Note that the symbol $|\cdot|$ in the left-hand side of the above formula denote groupoid cardinality: this is a weighted count of isomorphism classes, where each isomorphism class is weighted by the inverse of the size of its automorphism group. This is a better behaved notion of size for groupoids than naive isomorphism class count, since for example it interacts well with groupoid quotients. The precise form of the factors in the right-hand side is not very relevant for us, but they are in fact very explicit and easy to compute in practice.

We will reformulate the left-hand side of this formula. If we let $G=\ul{\Aut}(\Lambda)$, then we have an equivalence of groupoids
\[
\{\text{lattices of genus }\Lambda\} \isoto G(\Q)\backslash G(\A) /G(\A_0),
\] 
where $\A_0 = \R\times \prod_p \Z_p$ and $\A = \Q\otimes \A_0$ are the rings of integral adeles and adeles respectively. The idea is as follows: if $\Lambda, \Lambda'$ are two lattices of the same genus, then we can pick isomorphisms $\phi: \Lambda\otimes \A_0 \overset{\sim}{\to} \Lambda'\otimes \A_0$ and $\psi :\Lambda\otimes \Q \overset{\sim}{\to} \Lambda'\otimes \Q$ (the former directly by the definition of genus, and the latter by the Hasse-Minkowski theorem), and $\psi^{-1} \circ \phi$ defines an element of $G(\A)$. The quotient by $G(\A_0)$ on the right and $G(\Q)$ on the left exactly accounts for the ambiguity in the choice of $\phi$ and $\psi$.


If $\mu$ is a Haar measure on $G(\A)$, then
\[
\left| G(\Q)\backslash G(\A) /G(\A_0) \right| = \frac{\mu(G(\Q) \backslash G(\A))}{\mu(G(\A_0))}
\] 
(where both the numerator and denominator are finite as $G(\A_0)$ is compact and $G(\Q)$ is a cocompact discrete subgroup of $G(\A)$. Note that the right-hand side is clearly independent of the normalization of $\mu$, but, in fact, for any semisimple algebraic group $G$ over $\Q$, there is a canonical choice of normalization $\mu = \muT$, called the \defn{Tamagawa measure}. The denominator $\mu(G(\A_0))$ expands as a product $\mu(G(\R)) \cdot \prod_p \mu(G(\Z_p))$, and one can check that, for $G=\ul{Aut}(\Lambda)$, this essentially exactly matches the right-hand side of Siegel's mass formula\footnote{More precisely, they differ by some trivial factor coming from the fact that orthogonal groups are neither connected not simply-connected.} Therefore, the Siegel mass formula can be seen as a special case of the following conjecture of Weil for $G=\ul{Aut}(\Lambda)$:

\begin{conjecture}[Tamagawa number conjecture] \label{conj:tamagawa-number}
    Let $G$ be a connected simply-connected semisimple group over $\Q$. Then
    \[ \muT\left( G(\Q) \backslash G(\A) \right) = 1.\]
\end{conjecture}

\section{Function field analog}
We will now formulate a function field analog of Conjecture \ref{conj:tamagawa-number}. Let $X$ be a smooth projective curve over $\F_q$. We replace $\Q$ by the function field $K=K(X)$ and we replace the adeles by their analog
\begin{align*}
	\A_0 &= \prod_{x\in X}\OO_x \\
	\A &= \A_0 \otimes K.
\end{align*}
Let $G$ be a smooth group scheme over $X$ with connected fibers and a semisimple, simply connected generic fiber. In this case the analog of the double quotient considered earlier as a very natural interpretation, uniformly for any $G$: it parametrizes $G$-bundles on the curve $X$ (defined over $\F_q$). Namely, we have an equivalence of groupoids
\[
G(K)\backslash G(\A) /G(\A_0) \cong \Bun_G(\F_q),
\] 
where $\Bun_G$ denotes the moduli stack of $G$-bundle on the curve $X$. The idea behind this equivalence is that any $G$-bundle can be trivialized at the generic point and trivialized on a formal neighborhood of each point. The transition function between these two trivializations is an element of $G(\A_0)$, and the quotient on both sides accounts for the ambiguity in the choice of trivializations.

The analog of the product formula for $\left| G(\Q)\backslash G(\A) /G(\A_0) \right|$, or equivalently of Conjecture \ref{conj:tamagawa-number} in the function field setting turns out to be the following:

\begin{theorem}[Gaitsgory--Lurie] \label{thm:numerical-product-formula}
	\[
	\left| \Bun_G(\F_q) \right| = q^{\dim\Bun_G}\prod_{x\in X} \frac{|k_x|^{d}}{|G(k_x)|}
	\] 
	where $d=\dim G$.
\end{theorem}
The advantage of the function field setting is that the left-hand side is not just a groupoid, but it consists the $\F_q$ points of a nice stack, so we have more tools from algebraic geometry available, such \'etale cohomology. In particular, we have the Grothendieck-Lefschetz fixed point formula, which tells us that
\[
\left| \Bun_G(\F_q) \right| = \Tr(\Frob, H_c^{\bullet}(\Bun_G)) = q^{\dim \Bun_G}\Tr(\Frob^{-1}, H^{\bullet}(\Bun_G)).	
\] 
(where the second equality is Poincaré duality on the smooth stack $\Bun_G$). Theorem \ref{thm:numerical-product-formula} can therefore be reformulated as saying that 
\begin{equation} \label{eq:numerical-product-formula}
    \Tr(\Frob^{-1}, H^{\bullet}(\Bun_G)) = \prod_{x\in X} \frac{|k_x|^{d}}{|G(k_x)|}.
\end{equation}

\section{Cohomological product formula}

The numerical product formula (\ref{eq:numerical-product-formula}) can be deduced from a kind of product formula for the cohomology of $\Bun_G$ itself, which is the main result of \cite{abott}. In order to formulate this cohomological product formula, we first need to introduce the \defn{Ran Space} $\Ran(X)$. Informally speaking, $\Ran(X)$ is a prestack parametrizing non-empty finite subsets of $X$. One can construct on $\Ran(X)$ a sheaf $\Bred$, whose $!$-stalk at a subset $\{x_1,\dots,x_n\} \in \Ran(X)$ is given by
\[
\bigotimes_{i=1}^{n} H\red^{\bullet}\left( \frac{\text{pt}}{G_{x_i}} \right)
\] 
(where $G_x$ denotes the fiber of $G$ above a point $x\in X$).
The the cohomological product formula reads as follows:

\begin{theorem} \label{thm:coh-product-formula}
    Let $X$ be a smooth projective curve over an algebraically closed field $k$ and $G$ be a smooth group scheme over $X$ with connected fibers and a semisimple, simply connected generic fiber. There is a canonical isomorphism
    \[H_c^{\bullet}(\Ran,\Bred) \simto H\red^{\bullet}(\Bun_G).\]
\end{theorem}

Most of the rest of this seminar will be devoted to working through the proof of Theorem \ref{thm:coh-product-formula}. Fow the rest of this introduction, we will explain how to deduce the numerical product formula (\ref{eq:numerical-product-formula}) from it in the case where $k=\overline{\F}_q$ and $X$ and $G$ are defined over $\F_q$

\begin{remark}
    Informally, one can think of think of $\Bun_G$ as a sort of "continuous" product of $\frac{\text{pt}}{G_x}$ over all $x\in X$. Moreover, by the definition of the sheaf $\Bred$ the left-hand-side $H_c^{\bullet}(\Ran,\Bred)$ can be thought of as a sort of "continuous" tensor product of $H\red^{\bullet}\left( \frac{\text{pt}}{G_{x}}\right)$ over all $x\in X$. This is the reason behind the name "cohomological product formula".
\end{remark}

\section{Atiyah--Bott formula}

We will now combine the cohomological product formula with a classical explicit description of the cohomology of classifying spaces of reductive groups to deduce an explicit description of the cohomology of $\Bun_G$, called the Atiyah-Bott formula.

We first recall this classical description. Let $G_0$ be a semisimple group over an algebraically closed field $k$. Let $F=\mathbb{Q}_\ell$ be the field of coefficients of our cohomology theory. Then we have isomorphisms of graded rings
\[
	H^{\bullet}\left( \frac{\text{pt}}{G_0} \right) = F[\tf]^{W} \isoto F[f_1,\dots,f_r]
\] 
where $\tf$ is a Cartan subalgebra for $G_0$, $r=\dim \tf$ is the rank of $G_0$ and $\deg f_i = 2d_i$ (where $d_i$ are the exponents of $G_0$). Note that the second isomorphism is non-canonical.
\begin{example}
	Let $G_0 = \SL_n$, then we get
	\[
		H^\bullet\left(\frac{\text{pt}}{\SL_n} \right) = \frac{F[x_1,\dots,x_n]^{S_n}}{(x_1+\dots+x_n)} \isoto F[e_2,\dots,e_n],
	\] 
    where the $e_k$ are elementary symmetric polynomials. This is the fundamental theorem of symmetric polynomials.
\end{example}

This description can be reformulated as follows: let $\km$ be the augmentation ideal in $H^{\bullet}\left( \frac{\text{pt}}{G_0} \right) $ and
\[
M_0 = \frac{\km}{\km^2}.
\] 
Then we have non-canonical isomorphisms
\begin{equation} \label{eq:coh-pt/G}
	M_0 \cong \bigoplus_{i=1}^{r}F[-2d_i],\qquad H^{\bullet}\left( \frac{\text{pt}}{G_0} \right) \isoto \Sym^{\bullet}M_0.
\end{equation}
Moreover, if $k=\overline{\F}_q$ and $G_0$ is defined over $\F_q$, we can choose this isomorphism to be $\Frob$-equivariant and $\Frob$ acts on $F[-2d_i]$ by $q^{d_i}$.

We now come back to our global setting and let $G$ be a group scheme over a curve $X$. We will assume here to simplify the exposition that all fibers of $G$ are semisimple.  From $G$ we can construct a local system $M$ on $X$ whose fiber $M_x$ at $x\in X$ is the $M_0$ from the construction above applied to the fiber $G_0 = G_x$.

We are now ready to formulate the Atiyah-Bott formula and sketch its proof (using the cohomological product formula).

\begin{theorem}[Atiyah--Bott] \label{thm:atiyah-bott}
There is a non-canonical isomorphism
	\[
	H^{\bullet}(\Bun_G) \isoto \Sym^{\bullet}H^{\bullet}(X,M).
	\] 
If $k=\overline{\F}_q$ and $X,G$ are defined over $\F_q$, we can choose this isomorphism to be $\Frob$-equivariant.
\end{theorem}

\begin{proof}
	By Theorem \ref{thm:coh-product-formula}, we need to show that there is an isomorphism
	\[
	H_c^{\bullet}(\Ran,\Bred) \isoto \Sym^{+}\left( H^{\bullet}(X,M) \right),
	\] 
    where $\Sym^+$ denotes the positive degree part of the symmetric algebra.
    
	There exists a symmetric monoidal convolution product on $\Shv^{!}(\Ran)$ (the sheaf category in which $\Bred$ lives) such that $H_c^{\bullet}(-)$ is a symmetric monoidal functor, so it's enough to show that
	\begin{equation} \label{eq:bred-ab}
	    \Bred \isoto \Sym^{+}(i_{!}M)
	\end{equation}
	where $i:X\injto \Ran$ denotes the inclusion of $X$ as the locus of singletons in the Ran space.

	Both sides have the structure of a ``commutative factorization algebra'' in $\Shv^{!}(\Ran)$. This is a notion that we will study in more details later in the seminar, but roughly speaking a factorization algebra is an algebra object $\CF\in \Shv^{!}(\Ran)$ for which the algebra structure induces isomorphisms between the $!$-stalk at $\{x_1,\dots,x_n\} $ and the tensor product of $!$-stalks at each $x_i$. In particular these factorization algebras are essentially determined by their $!$-restriction to $X \overset{i}{\hookrightarrow} \Ran$.

	So it's enough to prove that both sides of (\ref{eq:bred-ab}) after $!$-restriction to $X\injto \Ran$. If $G=X\times G_0$ is constant, then the left-hand side is the constant sheaf with fiber $H\red^{\bullet}\left( \frac{\text{pt}}{G_0} \right)$ while the right-hand side is the constant sheaf with fiber $\Sym^{+}M_0$, so this follows immediatly from (\ref{eq:coh-pt/G}). In general, one needs to pass to a finite Galois cover $\widetilde{X} \to X$ over which $G$ becomes constant, and hence over which there exists an isomorphism. To show that there exists an isomorphism that descends back to $X$, one uses the following fact: over a field $F$ of characteristic zero, any affine space equipped with an action a finite group $\Gamma$ by affine transformations has a fixed point (in our case, the space of all possible isomorphisms forms an affine space acted on by the finite group $\Gamma = \operatorname{Gal}(\widetilde{X}/X)$).
\end{proof}

\section{Proof of the numerical product formula}
We will now deduce (\ref{eq:numerical-product-formula}) from Theorem \ref{thm:atiyah-bott}. We need to calculate
\begin{align*}
	\Tr(\Frob^{-1}, H^{\bullet}(Bun_G)) &= \Tr\left( \Frob^{-1}, \Sym^{\bullet}H^{\bullet}(X,M) \right)  \\
	&= \frac{1}{\det(1-\Frob^{-1},H^{\bullet}(X,M))}
\end{align*}
The right-hand side $\Tr\left( \Frob^{-1}, \Sym^{\bullet}H^{\bullet}(X,M) \right)$ is an infinite sum, which converges since the eigenvalues of $\Frob^{-1} $ on the fibers of $M$ are bounded by $q^{-4}$, hence its eigenvalues on $H^{\bullet}(X,M)$ are bounded by $q^{-2}$ by \cite{weilII}. Now, by the Grothendieck formula for $L$-functions \cite[Theorem 4.4]{freitag-kiehl}, we have:
\[
\frac{1}{\det\left( 1-t\Frob^{-1},H^{\bullet}(X,M) \right) } = \prod_{x\in X}\frac{1}{\det\left( 1-t^{d_x}\Frob_x^{-1}, M_x \right) }.
\] 
Moreover, note that
\begin{align*}
	\frac{1}{\det\left( 1-\Frob_x^{-1},M_x \right) } &= \Tr\left( \Frob^{-1}, \Sym^{\bullet}M_x \right)  \\
	&= \Tr\left( \Frob^{-1}, H^{\bullet}\left( \frac{\text{pt}}{G_x} \right)  \right)  \\
	&= \frac{|k_x|^{d}}{|G(k_x)|},
\end{align*}
where the second equality follows from (\ref{eq:coh-pt/G}) and the last one from the Grothendieck-Lefschetz trace formula applied to $\frac{\text{pt}}{G_x}$.

Putting all this together, we get
\begin{align*}
	q^{\dim\Bun_G}\left| \Bun_G(\F_q) \right| &= \Tr(\Frob^{-1}, H^{\bullet}(\Bun_G)) \\
						  &= \prod_{x\in X} \frac{|k_x|^{d}}{|G(k_x)|},
\end{align*}
which completes the proof of Theorem \ref{thm:numerical-product-formula}.

\section{Roadmap to the proof of the cohomological product formula}

Let's give a rough sketch of the strategy underlying the proof of Theorem \ref{thm:coh-product-formula}, which will be our main goal for the rest of this seminar. The key input (which will be taken as a black box) is non-abelian Poincar\'e duality (NPD). Consider the prestack $\Gr_{\Ran}$ which parametrizes triples of a point $S\in \Ran$, a $G$-bundle $\mathcal{P}\in \Bun_G$, and a trivilization of $\mathcal{P}$ away from $S$.
\begin{theorem}[Non-abelian Poincar\'e Duality]
	The projection map $\Gr_{\Ran}\to \Bun_G$ induces an isomorphism on homology.
\end{theorem}
We can reformulate NPD as follows. If $\CA$ denotes the $!$-pushforward of $\omega$ under $\Gr_{\Ran}\to \Ran$, then NPD is saying that the canonical map
\[
H_c^\bullet(\Ran,\CA) \to H_\bullet(\Bun_G).
\] 
is an isomorphism.
For slightly technical reasons, we will need to replace $\CA$ by a "reduced version" $\CA\red$, which satisfies
\[
H_c(\Ran,\CA\red) \simto H_\bullet ^{\text{red}}(\Bun_G).
\] 
Assuming NPD, we can therefore reduce the cohomological product formula to the statement that 
\[
	H_c^\bullet(\Ran,\CA\red) = H_c^\bullet(\Ran,\Bred)^{\vee}.
\] 
This will be a consequence of Verdier duality on the Ran space: namely, we will prove that $\Bred=\D\CA\red$, where $\D$ denotes Verdier duality.

There is indeed a canonical map $\varphi: \Bred\to \D\CA\red$, coming from a pairing between $\Bred$ and $\CA\red$. To show that it is an isomorphism, we will use again factorization algebra techniques: it will be shown that both $\CA\red$ and $\Bred$ have the structure of commutative factorization algebras, so it suffices to check that $\varphi$ is an isomorphism after $!$-restriction to  $X \overset{i}{\hookrightarrow} \Ran$. This will be checked $!$-stalkwise. The statement that $\varphi$ induces an isomorphism on $!$-stalks at a point $x\in X$ is called the \emph{local duality} statement. Two different proofs of this statement will be given (note that, since it is of local nature, it is enough to prove it for $G$ a constant group scheme and for any single choice of $X$, say $X=\P^1$).

\end{document}

